\section{Introduction}\label{sec:intro}

A manufacturer that supplies many markets rarely ships from the factory
gate to the customer in a single leg. Goods leave a small number of
production sites, pass through intermediate depots that consolidate and
break bulk, and only then reach the retailers. Deciding which factories
and which depots to operate, and how to route several products through
the two stages so that every demand is met at least total cost, is among
the oldest planning questions in distribution logistics
\citep{KloseDrexl2005,OrtizAstorquizaEtAl2018}, and it remains
a daily one wherever a supply chain is redesigned or a distribution
network is sized. The two layers interact from the start, because a depot
is worth opening only if some open factory can feed it and some customer
needs what it holds, so the fixed decisions and the transport decisions
cannot be taken apart, and it is this coupling that makes the problem more
than the sum of its two stages.

This interaction is captured by the multiproduct two-stage capacitated
facility location problem, which we abbreviate MP-TSCFLP. Several
products flow from candidate factories to candidate depots and on to
customers. Opening a factory or a depot costs a fixed charge, each unit
shipped on a stage carries a per-product transport cost, and every
factory and depot has a per-product capacity, while each customer has a
per-product demand that must be satisfied. Both stages are complete
bipartite, so any open factory may serve any depot and any open depot may
serve any customer. The uncapacitated single-product ancestor of the
model is the classical facility location problem, and the two-stage
capacitated multiproduct form was already observed to be \NP-hard by
\citet{PirkulJayaraman1998}. The benchmark and the solution methods we
build on are those of \citet{MauriEtAl2021}, who introduced the
hundred-instance benchmark and attacked it with two hybrid
metaheuristics, and of \citet{MoraisEtAl2022}, who later added local
branching to a genetic-algorithm hybrid for the same model.

Two facts recur in this literature, and both concern difficulty rather
than modelling. First, the instances resist exact solution. On the
hundred instances assembled by \citet{MauriEtAl2021} a commercial
mixed-integer solver solved none to proven optimality within a one-hour
budget per instance, which is why the authors turned to heuristics in the first place.
Second, this difficulty has no theoretical counterpart. Beyond the
inherited statement that the problem is \NP-hard
\citep{PirkulJayaraman1998}, no finer-grained account exists of \emph{where}
the hardness sits, which structural features drive it, or which of them could
be exploited to solve the problem faster.

This paper takes the empirical hardness as the phenomenon to be
explained. Facility location as such has a parameterized literature,
whose closest antecedent here is the parameterized view of
\citet{FellowsFernau2011}, a map of uncapacitated single-stage variants,
but no parameterized analysis covers the two-stage capacitated
multiproduct model, and we give what is, to our knowledge, the first. We
work throughout with the cardinality-constrained decision version, in which
a solution may open at most $k$ facilities, counting factories and depots
together, and $k$ is
the parameter that measures how lean a viable design is allowed to be.
Alongside $k$ we study the four size parameters of an instance, the numbers
$|I|$ of factories, $|J|$ of depots, $|K|$ of customers and $|L|$ of
products, both singly and in every combination. The guiding question is
which of these quantities, when held small, tames the exponential behaviour
of the problem, and which leave it intractable even when fixed. An
answer matters beyond taxonomy, because each tractable combination names
a regime in which exact solution scales, and each hard combination names
a feature whose cost no exact method avoids unless \Ppoly${}={}$\NP. The
analysis is a classification, so the algorithms it produces are devices for
locating the difficulty rather than candidates to beat a commercial solver,
and its practical residue is a covering bound that annotates an instance
before any optimisation is run.

Our findings fall into three groups, one locating the hardness along
the four dimensions, one drawing the tractability border together with
its preprocessing limits, and one certifying the map computationally and
confronting it with a solver. The hardness comes first because its
slices are the material from which the border is proved.

First, we locate the hardness. The problem collapses when all fixed
charges vanish and when a single factory faces a single depot
(Propositions~\ref{prop:fzero} and~\ref{prop:onebyone}), and hardness
begins one step above. With one product and one factory the problem is
already strongly \NP-hard, the depot-to-customer stage alone encoding a
covering problem, and no approximation beats the logarithmic Set Cover
threshold unless \Ppoly${}={}$\NP\ (Theorem~\ref{thm:cover},
Corollary~\ref{cor:coverinapprox}). Fixing any single size parameter to
a constant leaves the problem \NP-hard, which we prove through four
hard groups, each fixing two of the four dimensions to one. Behind the
groups sit four discrimination pairs, pairs in which one facility side
tells the members of another dimension apart through its cost matrix, a
notion Section~\ref{sec:complexity} fixes precisely, so hardness
survives every single restriction.

Second, we draw the tractability border. Parameterizing by any subset
$P$ of the four size parameters, the problem is fixed-parameter
tractable whenever $P$ contains both $|I|$ and $|J|$, through a single
$2^{|I|+|J|}\cdot\mathrm{poly}$ algorithm, and para-\NP-hard for every
other $P$, so under $\Ppoly\neq\NP$ the tractable subsets are exactly
those containing both (Theorem~\ref{thm:dichotomy}). MP-TSCFLP is
$\mathrm{W}[2]$-hard in the cardinality parameter $k$, already with one
product and one factory and already when only the depots are counted
(Theorem~\ref{thm:w2}), and under the Exponential Time Hypothesis the
naive $n^{k}$ enumeration is essentially optimal
(Theorem~\ref{thm:eth}). A separate numeric layer, our label for
hardness that persists after the combinatorial dimensions are
trivialised, governs the smallest instances. With a single customer and
a single product the problem is \NP-hard, separability of the
first-stage cost delimits a pseudo-polynomially solvable part, and
strong hardness persists outside the separable class, the general
boundary remaining open (Theorem~\ref{thm:numeric} and
Theorem~\ref{thm:separable}). On the preprocessing side, the problem
admits no polynomial kernel in $|I|+|J|$ unless
\NP${}\subseteq{}$coNP/poly (Theorem~\ref{thm:nokernel}), and over the
incidence representation we analyse the obstruction is numeric rather
than structural, for that incidence graph always has clique-width at
most two (Observation~\ref{obs:cliquewidth}). Several further positive
results, the pseudo-polynomial regimes and the
customer aggregation among them, are stated and classified but kept
outside the algorithmic core, since our aim is the map and not the
tuning of any one route through it.

Third, every construction behind the map is checked computationally.
The gadgets were validated by a shared battery of roughly
$9\times 10^{5}$ individual checks, none of which failed, and the enumeration algorithm agrees
with an independent brute-force solver across all $1{,}381$
comparisons. The computational study asks in what regime the
enumeration is practical and how it stands against Gurobi, the
mixed-integer baseline of Section~\ref{sec:experiments}. It then asks
whether the pruning threshold $k^{\ast}$ that our branch-and-bound
derives, available before any optimisation is run, tracks the
empirical difficulty of an instance, and through what mechanism it
does so. The protocol, instances and seeds are fixed in advance so
that every figure can be regenerated.

The remainder of this paper is organised as follows.
Section~\ref{sec:literature} places the model in the exact, heuristic and
parameterized literatures and closes with a positioning table.
Section~\ref{sec:prelim} fixes the model, proves the routing oracle,
that is the subroutine that, given a design, computes its optimal
routing cost, together with the
feasibility test and the monotonicity on which every later argument
rests, and maps the polynomial cases. Section~\ref{sec:complexity}
develops the four covering reductions, the numeric layer and the lower
bounds for the parameter $k$. Section~\ref{sec:algorithms} gives the
enumeration algorithm with its branch-and-bound, the dichotomy theorem
and the two sides of the kernelization question, and
Section~\ref{sec:experiments} reports the computational study built on
them. Section~\ref{sec:conclusion} states the two-layer synthesis and
five open problems.
